\begin{remark}[What this settles about the technique needed]
Only a one-sided bound together with an exceptional set of size $o(b)$ is
required. The precise asymptotic machinery of~\cite{DRZZ25} --- their strange
and pseudo-differentiable functions --- is therefore not needed here, and we
cite them for context only. The question was nevertheless the right one: the
non-smoothness of $\log|2\cos|$ is exactly where the elementary excision of
step~4 does its work.
\end{remark}

\begin{remark}[An alternative route]
One can instead avoid~\eqref{eq:fourier} altogether and argue through the full
discrepancy $D_B(\theta)$ of $\{p\theta/2\pi\}_{p\in B}$: bound it by
Erd\H{o}s--Tur\'an from the same exponential sums, then apply Koksma's
inequality to the truncated function $\max(\log|\cos\pi x|,\log\delta)$. The two
routes give the same conclusion. We prefer the one above for two reasons:
\eqref{eq:fourier} is an exact identity, and step~3 needs the Erd\H{o}s--Tur\'an
machinery only for a single interval, where it is elementary, rather than
uniformly over all intervals.
\end{remark}

\begin{remark}[Where the effective constants are, and are not]
Effectivity in this paper is confined to three places: the peak heights at
$q=4,6$, which are exact identities (Corollary~\ref{cor:spectrum}(i)--(ii)); the
remaining finite list of moduli, effective with $k_0\le1.5\cdot10^{4}$
(Proposition~\ref{prop:l5a}); and the Gaussian envelopes, which need no
arithmetic input at all (Proposition~\ref{prop:curv}). This section is
asymptotic; see Remark~\ref{rem:helfgott}.
\end{remark}

\begin{remark}[On explicit constants]\label{rem:helfgott}
One might hope to make step~3 effective using Helfgott's explicit minor-arc
bounds~\cite{Hel12}. His Main Theorem is stated for $x\ge x_0=2.16\cdot10^{20}$,
which corresponds to $k=\pi(x_0)\approx4.7\cdot10^{18}$. That is far outside any
range we could combine with the direct computations of Section~\ref{sec:num}, so
we do not claim an effective version of this section. Effectivity in this paper
is confined to Proposition~\ref{prop:l5a} (thresholds $\approx10^{4}$) and to
Proposition~\ref{prop:curv}, which needs no arithmetic input at all.
\end{remark}

\begin{remark}[Numerical check of the budget]\label{rem:budget}
The parameters are asymptotic and are meaningless at $b=23$, so we tested the
steps above at $b$ up to $1999$, taking the worst of six sampled deep-minor
$\theta$ at each size. Each evaluation is $O(bN_0)$. One substitution is forced:
the argument takes $N_0=(\log N)^{A'}$ with $A'=12$, which at $b=1999$ is about
$2\cdot10^{10}$ and cannot be computed; the table therefore uses
$N_0=(\log b)^{3}$, which is the same qualitative object --- a fixed power of a
logarithm --- and is the smallest such choice for which the tail is already
meaningful. Increasing the exponent only improves the tail and worsens the head,
both monotonically, so the comparison with $b^{1/4}$ below is not sensitive to
it. Writing $\mathrm{tail}$ for the step-4 bound $1/(N_0\delta_0)$:
\[
  \begin{array}{l|cccc}
    b & 99 & 299 & 999 & 1999\\\hline
    |B_{\mathrm{bad}}|/b\ \text{(measured)} & 0.485 & 0.405 & 0.324 & 0.282\\
    \text{step-3 bound}/b,\ N_0=(\log N)^{3} & 1.079 & 0.910 & 0.614 & 0.459\\
    b^{-1}\sum_{n\le N_0}|S_B(n\theta)|/n & 0.697 & 0.640 & 0.415 & 0.292\\\hline
    \mathrm{tail},\ N_0=b^{1/4}\ \text{(old)} & 1.532 & 1.425 & 1.381 & 1.267\\
    \mathrm{tail},\ N_0=(\log N)^{3}\ \text{(new)} & 0.047 & 0.031 & 0.021 & 0.017
  \end{array}
\]
The bound of step~3 holds in every case. The change of $N_0$ is the decisive
one: with $N_0=b^{1/4}$ the tail bound never drops below $1$ in any accessible
range --- it would need $b\approx10^{6}$ --- so that step was vacuous in
practice, whereas with $N_0=(\log N)^{3}$ it is already below $0.05$ at $b=99$.
The price is a larger head, since more harmonics are summed; the head is the
slower of the two now, and the step-3 bound only becomes nontrivial (that is,
smaller than $1$) from $b\approx300$.

None of this is close to binding. The conclusion needs only
$\tfrac12e^{o(1)}<\sqrt3/2$, i.e.\ the $o(1)$ to stay below
$\log\sqrt3=0.549$, and the quantity it controls,
$-b^{-1}\log|G_B(\theta)|-\log2$, was measured at $0.123$ for $b=99$ and $0.024$
for $b=1999$.
\end{remark}

\begin{remark}[What the peaks actually look like]\label{rem:peakscan}
As a check on the dissection we located every local maximum of $|G_B|$ on
$[0,\pi]$ by direct evaluation on a grid of $6\cdot10^{6}$ points, for
$k=40,60,100,140$. Every one of the twelve largest lies within the envelope
radius of a rational $2\pi j/q$ with $q\le26$, and the ordering by height
follows $M(q)$ as predicted, with one instructive exception: at $q=10$, where
$5\in B$ forces $G_B(2\pi/10)=0$ exactly, the vanishing cuts a narrow notch and
the flanking maxima sit at $b^{-1}\log|G_B|=-0.355$ at $b=99$, rising towards
$\log M(10)=-0.291$ like $-\tfrac32(\log b)/b$, which is the expected width of
the notch. Excluding the central peak and a disc of the envelope radius about
each rational with $q\le60$, the largest remaining value of
$b^{-1}\log|G_B|$ is $-0.305$ at $b=99$ and $-0.305$ at $b=139$, against
$\log(\sqrt3/2)=-0.144$: a margin of $0.16$ per element, stable in $k$.
\end{remark}
